\section{Convergence of the canonical Gaussian source drivers}
\label{S:source-driver-convergence}

We match the heat normalization, Laplace prefix and joint source law in
the author version~\cite{chen:author2026:time-dependency}*{(1.20), (1.21),
(2.9), (2.11), pp.~5--6, 8--9}. Given-noise retention, the Radon
implementation and the operator-path conclusions are proved here.
No first-moment variational theorem is an input.

Retain Hypothesis~\ref{H:homog}, the fixed $R,T,\zeta_R,E$ and parameters
of~\eqref{E:driver-parameters}, and the original noise space $H_W$.
Write $\mathscr B_\theta=C^\theta([0,T];\mathcal B(E))$.
Fix one raw test basis $f_j$ and its actual coefficient paths $A_j(t)$
from Proposition~\ref{P:actual-operator-driver}. Thus,
$M_N=\sum_{j\le N}W(f_j)A_j\to M_R$ in every finite $L^p(\mathscr B_\theta)$.
Zero covariance gives zero noises and drivers throughout; assume
$\Gamma\ne0$ in the construction.

\subsection{Gaussian transfer in the full path norm}

\begin{lemma}[Finite Gaussian linear images]\label{L:source-finite-comparison}
  Let $v_1,\ldots,v_N$ lie in a real Banach space and let $\xi$ be a
  centered Gaussian vector with covariance bounded by $CI$, $C\ge0$.
  For independent standard real $g_j$ and $p\ge1$,
  \begin{equation}\label{E:source-finite-comparison}
    \left(\E\left\|\sum_j\xi_jv_j\right\|^p\right)^{1/p}
    \le\sqrt C\left(\E\left\|\sum_jg_jv_j\right\|^p\right)^{1/p}.
  \end{equation}
\end{lemma}

\begin{proof}
  The function $\varphi(x)=\|\sum_jx_jv_j\|^p$ is continuous and convex,
  even when the linear image gives only a seminorm on $\R^N$.
  Its Gaussian expectations are finite by the finite-sum bound.
  On an auxiliary space take an independent centered Gaussian $Z$ of
  covariance $CI-\operatorname{Cov}\xi$. Conditional Jensen gives
  $\E\varphi(\xi)\le\E\varphi(\xi+Z)$. The latter vector has the law of
  $\sqrt C\,g$, proving the claim. The case $C=0$ is immediate.
  This auxiliary comparison changes neither the actual coupling of $\xi$
  nor any nonlinear wave solution.
  This proves Lemma~\ref{L:source-finite-comparison}.
\end{proof}

\begin{proposition}[Actual Gaussian transfer and stability]
  \label{P:source-gaussian-transfer}
  Let $Y$ be a centered Gaussian linear noise on smooth compact positive-time
  tests, on a common coupled space, with
  \begin{equation}\label{E:source-raw-variance-bound}
    \operatorname{Var}Y(F)\le C\|F\|_{H_W}^2.
  \end{equation}
  It extends to $H_W$. The finite paths
  $M_{Y,N}=\sum_{j\le N}Y(f_j)A_j$ converge in every finite
  $L^p(\mathscr B_\theta)$ to an actual centered Gaussian path $M_Y$ with
  compact values at all times on one path event, and
  \begin{align}
    \|M_Y-M_{Y,N}\|_{L^p\mathscr B_\theta}
    &\le\sqrt C\|M_R-M_N\|_{L^p\mathscr B_\theta},\label{E:source-transfer-tail}\\
    \|M_Y\|_{L^p\mathscr B_\theta}
    &\le\sqrt C\|M_R\|_{L^p\mathscr B_\theta}.\label{E:source-transfer-norm}
  \end{align}
  For a fixed real Schwartz pair and time,
  \begin{equation}\label{E:source-transferred-row}
    \langle M_Y(t)z,w\rangle_E=Y(h_{R,t;z,w})\quad\text{in }L^2,
  \end{equation}
  where $h_{R,t;z,w}$ is the literal raw test class
  $\zeta_R1_{(0,t)}v_zG_w$ with the actual adjoint from
  \eqref{E:variance-dual-wave}. The path is basis-independent up to
  indistinguishability, is linear in jointly Gaussian noise families, and
  complexifies sesquilinearly.

  If $(W,Y_n)$ are jointly centered Gaussian, $Y_n$ and $D_n=W-Y_n$
  have uniform bounds of the form~\eqref{E:source-raw-variance-bound},
  and $\operatorname{Var}D_n(F)\to0$ on every smooth compact test,
  then $M_{Y_n}\to M_R$ in every finite $L^p(\mathscr B_\theta)$.
\end{proposition}

\begin{proof}
  The variance bound annihilates raw-null tests and extends $Y$ as a
  bounded linear map into its actual $L^2$ space. Completion preserves
  Gaussian finite-dimensional laws. For each finite raw orthonormal family,
  the covariance of its $Y$ evaluations is bounded by $CI$, by applying
  the variance bound to every real linear combination. Independence of
  $Y(f_j)$ is not assumed.

  Apply Lemma~\ref{L:source-finite-comparison} with $v_j=A_j(\cdot)$ in
  $\mathscr B_\theta$, also to finite tails. The standard comparison vector
  has the law of the raw $W(f_j)$. The already-proved raw path convergence
  makes $M_{Y,N}$ Cauchy in every $L^p$. It lies in the same separable
  closed span of the countably many deterministic paths. The elementary
  Cauchy/subsequence construction of Lemma~\ref{L:driver-holder-limit}
  gives a strongly measurable limit on the actual space of $Y$.
  Different $p$ limits agree by convergence in probability. Passing to
  limits in the finite inequalities gives both displayed bounds. An
  almost surely uniformly convergent subsequence gives compact values
  for all times. Functional $L^2$ limits of Gaussian linear combinations
  prove that the path is Gaussian.

  For real Schwartz $z,w$, the actual coefficient identity is
  \[
    \langle A_j(t)z,w\rangle_E
    =\E[W(h_{R,t;z,w})W(f_j)]
    =\langle h_{R,t;z,w},f_j\rangle_{H_W}.
  \]
  Therefore, the finite transferred row is $Y(P_Nh)$, with $P_N$ the
  raw $H_W$ orthogonal projection. Since $P_Nh\to h$, the noise bound
  gives $Y(P_Nh)\to Y(h)$ in $L^2$. Operator-path $L^2$ convergence
  proves~\eqref{E:source-transferred-row}. Long intervals use the same
  raw completion and endpoint nullity as in Lemma~\ref{L:driver-interval-limit}.
  No filter is commuted through a physical mask or wave group. Four
  real pair identities give the complex assertion. Countably many dense
  Schwartz phases and rational times, followed by phase density and path
  continuity, prove basis indistinguishability.

  For jointly Gaussian $Y,Z$, their real linear combinations have a raw
  variance bound by the $L^2$ triangle inequality. The finite sums are
  linear in the noise, hence so are the path limits. In the stability
  assertion $D_n$ is centered Gaussian by the joint hypothesis. Its
  uniform bound gives
  \[
    \|M_{D_n}-M_{D_n,N}\|_{L^p\mathscr B_\theta}
    \le\sqrt{C_D}\|M_R-M_N\|_{L^p\mathscr B_\theta}.
  \]
  For fixed $N$, the covariance matrix of $(D_n(f_j))_{j\le N}$ is
  positive and bounded by $s_N(n)I$, where
  $s_N(n)=\sum_{j\le N}\operatorname{Var}D_n(f_j)\to0$.
  The finite comparison gives
  $\|M_{D_n,N}\|_{L^p\mathscr B_\theta}
  \le\sqrt{s_N(n)}\|M_N\|_{L^p\mathscr B_\theta}$.
  Choose $N$ large for the uniform tail, then $n$ large for the finite
  head. Linearity identifies $M_{D_n}=M_R-M_{Y_n}$. The same argument
  proves a multiparameter limit whenever each fixed test error tends
  to zero, with no relation between parameter rates.
  This proves Proposition~\ref{P:source-gaussian-transfer}.
\end{proof}

\subsection{The spatial covariance space and a retaining extension}

\begin{lemma}[Original spatial space and heat]\label{L:source-spatial-heat}
  Let $H_\Gamma$ be the real test-null completion with
  \[
    B_\Gamma(f,g)=\int\Gamma(\ud y)\int f(x)g(x-y)\ud x
    =(2\pi)^d\int\widehat f(\xi)\overline{\widehat g(\xi)}\mu(\ud\xi).
  \]
  It is separable and includes Schwartz tests. With
  \begin{equation}\label{E:source-heat-definition}
    p_\varepsilon(x)=(2\pi\varepsilon)^{-d/2}
    e^{-|x|^2/(2\varepsilon)},\qquad C_\varepsilon f=p_\varepsilon*f,
  \end{equation}
  heat is a positive self-adjoint contraction on $H_\Gamma$,
  $I-C_\varepsilon$ is a contraction, and $C_\varepsilon\to I$ strongly.
\end{lemma}

\begin{proof}
  The spectral normalization is the spatial computation in
  Proposition~\ref{P:masked-noise-completion}; its integral is real on
  real tests by evenness of $\mu$. The local bound
  $|B_\Gamma(f,g)|\le |K|\Gamma(K-K)\|f\|_\infty\|g\|_\infty$
  on a common compact support gives separability by countable smooth
  approximation. Expanding smooth cutoffs of a Schwartz test converge
  in Schwartz seminorms, and its covariance correlation pairing converges
  against tempered $\Gamma$. Thus, Schwartz tests enter the same completion.
  The Fourier multiplier of $C_\varepsilon$ is
  $m_\varepsilon(\xi)=e^{-\varepsilon|\xi|^2/2}$.
  On dense Schwartz tests, the displayed norm gives positivity,
  self-adjointness and both contraction bounds. Dominated convergence
  gives strong convergence there; uniform boundedness extends it to the
  completion. This uses no physical mask operator on a noise quotient.
  This proves Lemma~\ref{L:source-spatial-heat}.
\end{proof}

Set $\nu_a(\ud\lambda)=\lambda^{a-1}\ud\lambda/\Gamma_E(a)$ and
\[
  \mathcal G=L^2((0,\infty)_\lambda\times\R_r,
  \nu_a(\ud\lambda)\ud r;H_\Gamma).
\]
This is a real separable Hilbert space: countable simple functions on
rational rectangles with a countable dense spatial collection suffice.

\begin{lemma}[Feature isometry retaining the given noise]\label{L:source-retaining-isometry}
  On finite smooth temporal-spatial tensors define
  \begin{equation}\label{E:source-feature-isometry}
    \Phi F(\lambda,r)=\sqrt{2\lambda}
    \int_{\substack{t>0\\t\ge r}}e^{-\lambda(t-r)}F(t,\cdot)\ud t.
  \end{equation}
  It extends to an isometry $\Phi:H_W\to\mathcal G$ with closed range,
  and the formula represents smooth compact tests almost everywhere.
  On a product extension retaining the original $W$, there is an isonormal
  $\mathcal I$ on $\mathcal G$ with
  \begin{equation}\label{E:source-given-W-retention}
    \mathcal I(\Phi h)=W(h)\quad\text{in }L^2
    \quad\text{for each }h\in H_W.
  \end{equation}
\end{lemma}

\begin{proof}
  The identity
  \[
    \int_{-\infty}^{\min(t,u)}
    2\lambda e^{-\lambda(t+u-2r)}\ud r
    =e^{-\lambda|t-u|}
  \]
  and the raw Laplace formula give
  $\|\Phi F\|_{\mathcal G}^2=\|F\|_{H_W}^2$.
  All signed exchanges for finite tensors are absolutely integrable:
  compact temporal factors give integrals of
  $|\tau(t)\sigma(u)||t-u|^{-a}$, finite for $a<1$, times fixed finite
  spatial inner products. No oscillatory infinite-measure Fubini is used.
  Smooth time partitions on one compact support approximate a general
  smooth spacetime test uniformly by finite tensors. The local covariance
  bound~\eqref{E:driver-test-continuity} proves their density in $H_W$.
  Hence, the isometry extends with closed range.

  For a general smooth compact $F$, those approximants also converge
  uniformly in spatial $H_\Gamma$ norm in time, by local Radon mass.
  For fixed $\lambda,r$, their integrals in~\eqref{E:source-feature-isometry}
  converge in $H_\Gamma$. Meanwhile the isometry gives convergence in
  $\mathcal G$, with an almost everywhere convergent subsequence. Its
  two limits agree, proving the stated formula. The coordinate $r$ is
  auxiliary; no value of raw negative-time noise has been assigned.

  Let $P$ project orthogonally onto $\Phi H_W$. On a product extension
  take an independent isonormal $J$ on $(\Phi H_W)^\perp$, built from
  independent real standard Gaussians on a finite or countable basis.
  Define
  \[
    \mathcal I(g)=W(\Phi^{-1}Pg)+J((I-P)g).
  \]
  Orthogonality and independence make it isonormal on $\mathcal G$ and
  prove~\eqref{E:source-given-W-retention} with the actual original random
  variable, not only its law. For the countable fixed raw basis, these
  equalities hold on one common probability-one set. Thus, the raw
  finite paths and their $L^p$ limits are precisely pullbacks of the
  already constructed raw driver. Cutoff fields may use $J$; they are
  not asserted to be functions of $W$ alone.
  This proves Lemma~\ref{L:source-retaining-isometry}.
\end{proof}

\subsection{Exact joint source law and convergence}

Let $P_\eta$ multiply $\mathcal G$ by $1_{\{\lambda\le\eta^{-1}\}}$.
Let $C_\varepsilon$ act fiberwise on $H_\Gamma$. These commute. For
$\rho=(\varepsilon,\eta)>0$, define
\begin{equation}\label{E:source-cutoff-noises}
  Q_\rho=C_\varepsilon P_\eta,\qquad
  Y_\rho(F)=\mathcal I(Q_\rho\Phi F),\qquad
  D_\rho(F)=\mathcal I((I-Q_\rho)\Phi F)=W(F)-Y_\rho(F).
\end{equation}
The last equality uses the actual retaining coupling. Define
\begin{equation}\label{E:source-covariance-functions}
  \Gamma_e(x)=(p_e*\Gamma)(x),\qquad
  K_\eta(v)=\frac1{\Gamma_E(a)}
  \int_0^{\eta^{-1}}e^{-\lambda|v|}\lambda^{a-1}\ud\lambda.
\end{equation}
The temporal operation is a Laplace prefix, not temporal convolution.

\begin{proposition}[Canonical sources on the retaining space]
  \label{P:source-canonical-convergence}
  The noises in~\eqref{E:source-cutoff-noises} have pointfield realizations
  with joint covariance
  \begin{equation}\label{E:source-joint-covariance}
    \operatorname{Cov}(Y_{\varepsilon,\eta}(t,x),
    Y_{\varepsilon',\eta'}(s,y))
    =K_{\max(\eta,\eta')}(t-s)\Gamma_{\varepsilon+\varepsilon'}(x-y).
  \end{equation}
  Mean-square pointfield integration agrees with $Y_\rho(F)$ on tests.
  Both source and full residual obey
  \begin{equation}\label{E:source-two-dominations}
    \operatorname{Var}Y_\rho(F),\quad\operatorname{Var}D_\rho(F)
    \le\|F\|_{H_W}^2,
  \end{equation}
  and $\operatorname{Var}D_\rho(F)\to0$ for every $F\in H_W$ as
  $\varepsilon,\eta\downarrow0$ at arbitrary relative rates. Consequently,
  for every finite $p\ge1$,
  \begin{equation}\label{E:source-driver-limit}
    M_R^\rho:=M_{Y_\rho}\longrightarrow M_R
    \quad\text{in }L^p(\mathscr B_\theta),\qquad
    M_{D_\rho}=M_R-M_R^\rho.
  \end{equation}
  The source and residual have uniform local Gaussian driver moments and
  square-exponential norms, with constants for fixed $R,T$ and exponents.
\end{proposition}

\begin{proof}
  Apply the spatial distribution identity to a single shifted heat test:
  \[
    \Gamma_e(x)=\int e^{i\xi\cdot x}e^{-e|\xi|^2/2}\mu(\ud\xi).
  \]
  Gaussian spectral moments of every order are finite for tempered $\mu$.
  Thus, this is smooth with bounded derivatives. Since
  $\widehat{p_\varepsilon(x-\cdot)}
  =c_de^{-i\xi\cdot x}e^{-\varepsilon|\xi|^2/2}$, the two heat factors
  in the spatial inner product cancel $(2\pi)^d$ exactly, giving
  \begin{equation}\label{E:source-heat-inner-product}
    B_\Gamma(p_\varepsilon(x-\cdot),p_{\varepsilon'}(y-\cdot))
    =\Gamma_{\varepsilon+\varepsilon'}(x-y).
  \end{equation}
  Evenness of $\mu$ removes the reflected argument.

  For each positive $\rho$ and point $(t,x)$ put
  \begin{equation}\label{E:source-point-feature}
    v_{\rho;t,x}(\lambda,r)
    =1_{\{\lambda\le\eta^{-1}\}}\sqrt{2\lambda}
    e^{-\lambda(t-r)}1_{\{r\le t\}}p_\varepsilon(x-\cdot).
  \end{equation}
  Its squared $\mathcal G$ norm is
  $K_\eta(0)\Gamma_{2\varepsilon}(0)<\infty$.
  Define $Y_\rho(t,x)=\mathcal I(v_{\rho;t,x})$. The common $\lambda$
  range for two such vectors is
  $\lambda\le\min(\eta^{-1},\eta'^{-1})=\max(\eta,\eta')^{-1}$.
  Integrating $r$ gives $e^{-\lambda|t-s|}$, and
  \eqref{E:source-heat-inner-product} gives the spatial factor. This
  proves~\eqref{E:source-joint-covariance} with the prescribed widths and
  raw amplitudes, also for singular Radon and white covariances.

  The point feature is continuous in $\mathcal G$ norm, directly from
  this continuous covariance. Its norm is bounded on compact parameter
  sets. Hence, for a smooth compact spacetime $F$, its Bochner integral
  against $v_{\rho;t,x}$ exists. Integrate the heat translates in $x$,
  then the time feature, to obtain exactly $Q_\rho\Phi F$. Applying the
  bounded map $\mathcal I:\mathcal G\to L^2$ proves the pointfield-test
  identity on the same space. Common-support indicator tests follow by
  raw completion and $L^1$ approximation; bounded point-feature norms and
  the domination below give the same $L^2$ limit.

  Orthogonality of the two $\lambda$ ranges gives, for $g\in\mathcal G$,
  \[
    \|(I-Q_\rho)g\|_{\mathcal G}^2
    =\|(I-P_\eta)g\|_{\mathcal G}^2
    +\|P_\eta(I-C_\varepsilon)g\|_{\mathcal G}^2
    \le\|g\|_{\mathcal G}^2.
  \]
  Also $\|Q_\rho g\|_{\mathcal G}\le\|g\|_{\mathcal G}$. Together with
  the isometry these prove both variance bounds. Disjoint $\lambda$
  ranges are independent; the full heat residual need not be independent
  of the source.

  The $L^2$ tail integral gives $P_\eta\to I$ strongly. Fiberwise heat
  converges strongly on simple $\mathcal G$ functions by
  Lemma~\ref{L:source-spatial-heat}, then on all of $\mathcal G$ by the
  contraction bound. The estimate
  \[
    \|(I-Q_\rho)g\|_{\mathcal G}
    \le\|(I-C_\varepsilon)g\|_{\mathcal G}
    +\|(I-P_\eta)g\|_{\mathcal G}\longrightarrow0
  \]
  holds at arbitrary relative rates. Thus, the residual variance tends
  to zero for every raw $F$. Proposition~\ref{P:source-gaussian-transfer}
  gives~\eqref{E:source-driver-limit} without commuting a filter through
  $\zeta_R$ or a wave group.

  The same comparison gives
  $\|M_R^\rho\|_{L^p\mathscr B_\theta},
  \|M_{D_\rho}\|_{L^p\mathscr B_\theta}
  \le\|M_R\|_{L^p\mathscr B_\theta}$. The raw $\sqrt p$ estimate and
  its even-moment exponential-series argument give uniform Gaussian
  driver bounds. They are statements about $L^p$ equivalence classes
  on the retaining extension, not a simultaneous path event for all
  uncountably many cutoffs.
  This proves Proposition~\ref{P:source-canonical-convergence}.
\end{proof}

\subsection{Regularized pointfields and the ordinary first response}

\begin{lemma}[Regularity for each fixed cutoff]\label{L:source-pointfield-regularity}
  For each fixed $\rho>0$, the pointfield has a version continuous in
  time and smooth in space, with every spatial derivative jointly
  continuous on compact time-space rectangles.
\end{lemma}

\begin{proof}
  For each spatial multiindex $b$, differentiate
  \eqref{E:source-point-feature} in $x$ as a $\mathcal G$ vector.
  All derivatives exist because their spatial squared norms are the
  finite heat moments
  \[
    \int \xi^{2b}e^{-\varepsilon|\xi|^2}\mu(\ud\xi),
  \]
  and the next moments dominate difference quotients. Set
  $Z_b(t,x)=\mathcal I(\partial_x^bv_{\rho;t,x})$.
  At fixed $x$, use $1-e^{-\lambda h}\le\lambda h$ in the finite
  $\lambda$ integral; at fixed $t$, use
  $|e^{i\xi\cdot x}-e^{i\xi\cdot y}|\le|\xi||x-y|$ and the next heat
  moment. The squared triangle inequality gives
  \begin{equation}\label{E:source-derivative-increments}
    \E|Z_b(t,x)-Z_b(s,y)|^2
    \le C_{\rho,b}(|t-s|+|x-y|^2).
  \end{equation}
  This computation uses~\eqref{E:source-joint-covariance}, including its
  spatial width, and applies to singular spectral measures without a
  density assumption.

  Here is a direct continuity construction. In $m=d+1$ parameter
  dimensions on a compact rectangle, even Gaussian $q$th increment
  moments have root bounded by $C\sqrt q$ times the parameter distance
  to the power $1/2$ at small distances. The dyadic grid has at most
  $C_m2^{nm}$ adjacent edges at level $n$. Their maximum has $L^q$
  norm at most $C\sqrt q\,2^{-n(1/2-m/q)}$, with constants for this
  rectangle and fixed $\rho,b$. Choose $q$ large and
  $0<h_0<1/2-m/q$. The sum of these maxima weighted by $2^{nh_0}$
  has finite $L^q$ norm. Binary refinement, using a fixed number of
  adjacent edges per coordinate at each scale, gives a continuous
  extension from dyadic points. At each fixed parameter the extension
  agrees with the given Gaussian variable in $L^2$, by
  \eqref{E:source-derivative-increments}; it is a modification.
  Use countably many nested compact rectangles and all multiindices.
  The versions agree on overlaps on countable dense sets, hence
  everywhere there by continuity, so they patch.

  They are actual derivatives. The $\mathcal G$-valued fundamental
  theorem gives, along a coordinate segment,
  \[
    \partial_x^bv(t,x+he_i)-\partial_x^bv(t,x)
    =\int_0^h\partial_x^{b+e_i}v(t,x+ue_i)\ud u.
  \]
  Apply bounded linear $\mathcal I$ and $L^2$ integration. The continuous
  versions give the same ordinary integral by their Riemann sums, which
  converge both pathwise and in $L^2$. Enforce these identities first on
  countably many rational $t,x,h$ and multiindices, then use joint
  continuity to extend them to every segment in the compact rectangles.
  The ordinary fundamental theorem yields $\partial_iZ_b=Z_{b+e_i}$.
  Thus, all spatial derivatives exist and are jointly continuous with
  time. The versions are for this fixed $\rho$, not a common path event
  over uncountably many source parameters.
  This proves Lemma~\ref{L:source-pointfield-regularity}.
\end{proof}

\begin{proposition}[Actual ordinary source response]\label{P:source-ordinary-response}
  For each fixed $\rho>0$, on its regularity event, put
  $b_\rho(t,x)=\zeta_R(x)Y_\rho(t,x)$. The statewise integral
  \begin{equation}\label{E:source-ordinary-response}
    L_\rho(t)z=\int_0^tS(-u)B(b_\rho(u))S(u)z\ud u
  \end{equation}
  defines an operator-norm Lipschitz path on $E$, and
  $L_\rho=M_R^\rho$ for all $t\in[0,T]$ on one probability-one set.
  In particular, it has the strongly measurable path version supplied
  by the Gaussian construction.
\end{proposition}

\begin{proof}
  The coefficient is compact smooth in space and continuous in time in
  every spatial derivative norm. Multiplication estimates and strong
  continuity of the full free group show that the statewise integral
  exists on every $E_r$. Its operator increments $E_r\to E_{r+1}$
  are bounded by a finite sample constant times $|t-s|$, exactly by the
  statewise gain argument of Lemma~\ref{L:driver-gain-localization}.
  This proves Lipschitz continuity in $\mathcal B(E)$, without requiring
  norm continuity or an operator-norm derivative of $S(t)$.

  At each fixed $u$, the Sobolev duality calculation uses only spatial
  smoothness of $b_\rho(u,\cdot)$. From $S(-u)(0,f)$ and the two Bessel
  weights, for real Schwartz phases,
  \[
    \langle S(-u)B(b_\rho(u))S(u)z,w\rangle_E
    =\int b_\rho(u,x)[S(u)z]_1(x)G_w(u,x)\ud x.
  \]
  Thus, no prior hypothesis of smoothness in time is invoked for this
  instantaneous identity. For complex pairs conjugate the second factor.
  Integrating gives the ordinary source-field integral against the literal
  raw wave test $h_{R,t;z,w}$. For smooth tests, the point-feature integral
  has already been identified with $Y_\rho(F)$. The continuous pointfield
  version has the same integral: compact Riemann sums converge pathwise
  and in $L^2$ to the ordinary and feature integrals, respectively.
  Approximate the indicator wave test in the raw completion with common
  compact spatial support and $L^2$-time spatial supremum error tending
  to zero. Its $L^1$ error also tends to zero. The point-feature norms
  are bounded on that support, so their integrals converge in $\mathcal G$;
  the raw variance domination gives the same $L^2$ limit on the test side.
  Consequently,
  \[
    \langle L_\rho(t)z,w\rangle_E
    =Y_\rho(h_{R,t;z,w})
    =\langle M_R^\rho(t)z,w\rangle_E
    \quad\text{in }L^2.
  \]
  Compare a countable dense Schwartz phase collection and rational times.
  Boundedness and the path continuity of both sides extend this equality
  to all states and all times on one probability-one set. The equality
  is on the actual coupled space, not only in law. The set is for fixed
  $\rho$; the $L^p$ limit in~\eqref{E:source-driver-limit} needs no
  simultaneous version for every source parameter.
  This proves Proposition~\ref{P:source-ordinary-response}.
\end{proof}

The source laws above match the pinned author definitions. The author's
page~9 redefinition keeps the law; the retaining extension was proved in
Lemma~\ref{L:source-retaining-isometry}. The derivative estimate here was
derived from the correct joint covariance, not from the differing spatial
index in the author's page~8 increment display. No source correction is
adopted. All limits are at fixed finite $R,T$. The nonlinear Young flow,
canonical initial datum $(1,0)$ and its solution-moment transfer remain
separate parts of the argument.
